%%%%%%%%%%%%%%%%%%%%%%%%%%%%%%%%%%%%%%%%%%%%%%%%%%%%%%%%%%%%%%%%%%%%%%%%%%%%%%%%
% Official RA-L initial-submission style, Chinese advisor-facing outline.
% This file keeps the PaperCept/RA-L ieeeconf layout, but uses ctex for Chinese.
% Compile with XeLaTeX. This is for advisor discussion, not direct submission.
%%%%%%%%%%%%%%%%%%%%%%%%%%%%%%%%%%%%%%%%%%%%%%%%%%%%%%%%%%%%%%%%%%%%%%%%%%%%%%%%

\documentclass[letterpaper, 10 pt, conference]{ieeeconf}

\IEEEoverridecommandlockouts
\overrideIEEEmargins

\usepackage{amsmath,amssymb}
\usepackage{graphicx}
\usepackage{booktabs}
\usepackage{xcolor}
\usepackage[UTF8]{ctex}

\newcommand{\method}{ToolFrame}
\newcommand{\todo}[1]{\textcolor{red}{TODO: #1}}
\newcommand{\claim}[1]{\textbf{核心主张：}#1}
\newcommand{\openitem}[1]{\textit{待讨论：}#1}

\title{\LARGE \bf
\method{}：让工具定义灵巧操作的功能坐标系
}

% RA-L 初投稿为双盲；导师讨论版不填写作者信息。
\author{}

\begin{document}

\maketitle
\thispagestyle{empty}
\pagestyle{empty}

\begin{abstract}
这份大纲把 \method{} 的论文主线建立在一个底层机理上：工具使用的成败不由工具类别本身决定，而由工具的功能几何是否与场景中的目标物理效果对齐决定。真正困难的是，这个机理并不会直接出现在语言指令、物体类别或 VLM 生成的轨迹中；它必须从公开视觉证据中恢复出来，被表达成可计算的几何约束，并最终由灵巧手策略执行。本文拟将 \method{} 包装为一个工具为中心的 image-grounded tool-use pipeline：VLM 只负责选择稀疏、公开、可审计的工具和场景证据；几何编译器负责把工具接触、工具运动方向与场景效果绑定为 whole-tool $SE(3)$ 子目标；SimToolReal fixed-goal policy 负责执行；task contracts 负责最终成功率和失败归因。核心卖点是：本文不是让 VLM 生成机器人轨迹，而是把工具使用的隐藏功能机理显式化、几何化、可执行化。
\todo{最终摘要只在 launcher summary 和 task contract 冻结后加入实验数字。}
\end{abstract}

\section{INTRODUCTION}

\subsection{底层机理：工具不是类别，而是功能几何}
论文开头应从物理机理讲起，而不是从 VLM 规划讲起。工具之所以有用，不是因为它属于某个语义类别，而是因为它的某个几何部位能够把接触、力或运动传递到场景中。刷子通过刷毛产生扫动效果，橡皮通过平面接触产生擦除效果，锤子通过锤面产生冲击效果，螺丝刀通过尖端和杆轴产生插入/旋入效果，马克笔通过笔尖产生书写轨迹。对于灵巧手操作来说，手可以用很多方式抓住同一个工具，但任务成功只发生在少数 whole-tool 位姿中：功能部位接触正确，功能方向也与场景效果方向一致。

\subsection{核心困难：工具机理、场景效果和执行接口的绑定}
本文的底层挑战不是“识别工具”，而是一个三方绑定问题：
\begin{equation}
  \begin{aligned}
  &\text{工具功能机理}\\
  &\quad \longleftrightarrow \text{场景目标效果}
  \longleftrightarrow \text{整工具执行接口}.
  \end{aligned}
\end{equation}
许多 VLM-centric 机器人方法会把这个绑定过程隐藏在 VLM 输出的动作、稠密 waypoint 或轨迹中。这样工程上看起来直接，但科学上不够干净：成功可能来自隐藏目标先验、任务族 fallback、控制器行为或 prompt shortcut，而不一定来自真正的 image-grounded tool-use reasoning。

\subsection{核心洞见：工具使用更像编译问题，而不是轨迹生成问题}
本文最需要包装的洞见是：工具使用应被表述为一个 compiler problem，而不是 VLM trajectory-generation problem。VLM 的角色是选择稀疏公开证据，例如哪个工具部位有功能、场景里哪些点和方向定义了任务效果；几何编译器的角色是把这些证据转成整工具的度量约束；策略的角色只是执行编译出的固定目标。这样，整条因果链是可见的：
\begin{quote}
公开证据 $\rightarrow$ 功能 refs $\rightarrow$ whole-tool $SE(3)$ subgoals
$\rightarrow$ 灵巧手 rollout $\rightarrow$ task contract.
\end{quote}

\subsection{论文定位}
\claim{\method{} 是一个 image-grounded、tool-centric 的灵巧工具使用 pipeline：工具定义操作坐标系，编译器把它变成可执行 whole-tool 子目标，灵巧手负责跟随。}

这篇文章应避免被理解成普通的“VLM 机器人规划”或“工具 affordance 预测”。它不是让 VLM 写轨迹，因为 VLM 只选择公开稀疏证据；它也不只是预测 affordance，因为工具 refs 会进入几何编译器，被转化成可执行的 $SE(3)$ 子目标。真正创新点是把 public tool evidence、public RGB-D scene evidence 和 fixed-goal dexterous execution 之间的桥搭起来。

\subsection{拟强调贡献}
\begin{enumerate}
  \item \textbf{机理显式化的工具使用表述。} 将工具使用表述为工具侧功能机理与场景侧目标效果的对齐问题。
  \item \textbf{受证据边界约束的 VLM 角色。} VLM 只选择稀疏的公开工具和场景 refs，不生成轨迹、不生成 primitive、不使用 task-family fallback。
  \item \textbf{工具为中心的几何编译器。} 将 contact point 和 motion axis 与 scene refs 编译为 whole-tool $SE(3)$ 子目标。
  \item \textbf{task-contract 评估协议。} 不只报告最终成功率，还沿着机理链条归因失败位置。
\end{enumerate}

\section{RELATED WORK POSITIONING}

\subsection{VLM for Robot Planning}
相关工作可以承认 VLM 是强语义和视觉接口，但指出它并不天然保证接触几何和工具功能方向正确。我们的区别要讲清楚：\method{} 不把 VLM 当 trajectory generator，而把它当 evidence selector；度量几何和执行可行性由显式模块负责。

\subsection{Tool Affordance and Functional Part Grounding}
affordance 和 functional part 相关工作通常能指出工具哪里可能有用，但它们往往停留在语义或局部区域层面。本文要强调缺口在于：functional part label 还不是 robot-executable constraint。我们的 contact point 和 motion axis 是会被 compiler 消费的几何接口。

\subsection{Constraint-Based Manipulation}
约束式操作提供了正确的几何语言，但通常默认 task frame、接触点或目标轴已经给定。本文的价值是从公开图像证据中产生这些约束，并且以 whole-tool 而不是末端执行器轨迹为核心。

\subsection{Tool-Use Benchmark}
benchmark 部分的贡献不应只说“我们有很多任务”，而应说“我们能解释失败发生在哪里”。这让失败 rollout 也成为有价值的诊断证据，而不是简单负样本。

\section{PROBLEM STATEMENT}

\subsection{紧凑技术接口}
导师版大纲中技术接口保持简洁即可：
\begin{equation}
  \mathcal{R}_{tool} =
  \{p^{contact}_{tool}, a^{motion}_{tool}\},
\end{equation}
\begin{equation}
  \mathcal{R}_{scene} =
  \{p^{start}_{scene}, p^{path}_{scene}, p^{end}_{scene},
  a^{direction}_{scene}\}.
\end{equation}
论文的问题是：如何从公开证据恢复这些 refs，并把它们编译成 fixed-goal policy 可执行的 whole-tool poses $T_{world,tool}\in SE(3)$。

\subsection{证据边界}
证据边界本身就是论文可信度的一部分。允许输入包括：public RGB-D scene evidence、public tool-part evidence、tool mesh、task language，以及执行所需的机器人状态。不允许使用 hidden evaluator target refs、oracle target axes、oracle trajectories、primitive fallback、repair loop、task-family fallback 或 seed-specific shortcut。

\subsection{为什么边界重要}
如果方法可以 snap 到 simulator target，或者在失败时退回手写 primitive，那么 image-grounded tool use 的核心主张就不可检验。本文应被包装成更严格、更可审计的路线：用更少的特权信息，换取更清晰的中间机理和失败归因。

\section{METHOD OUTLINE}

\subsection{总览图应表达什么}
Fig.~\ref{fig:overview} 应该是主卖点图。它不应只是软件流程图，而应呈现“机理分解”：工具侧机理、场景侧效果、编译器桥接、灵巧手执行。

\begin{figure*}[t]
\centering
\fbox{\parbox{0.92\textwidth}{\centering
中文导师版 Fig. 1 占位：
(A) 恢复工具功能机理；
(B) grounding 场景效果 refs；
(C) 编译 whole-tool $SE(3)$ subgoals；
(D) fixed-goal dexterous rollout + task-contract evaluation。}}
\caption{\method{} 概念图。公开工具证据定义 contact 和 motion refs；公开 RGB-D 场景证据定义 effect refs；几何编译器将二者绑定为 whole-tool subgoals；fixed-goal 灵巧手策略执行。}
\label{fig:overview}
\end{figure*}

\subsection{A. Tool-Centric Prior Extraction}
\claim{恢复工具的功能机理。}

这一节应解释为什么 object category 和 visible geometry 不够。方法输出应被概括为一个紧凑工具侧表示：functional part、contact point 和 motion axis。暂时不要写死尚未冻结的实现细节。重点是：输出不是语义标签，而是后续编译器使用的几何机理。

\subsection{B. Scene Key-Points Grounding}
\claim{恢复场景中的稀疏目标效果机理。}

这一节应强调任务不需要 VLM 画机器人轨迹，而需要 VLM 从公开 RGB-D evidence 中选择定义任务效果的稀疏 refs：起点、路径点、终点和方向。RGB-D lifting 将图像 refs 转为 3D metric refs。prompt、sampling、threshold 等未冻结细节暂时不放进导师版主线。

\subsection{C. Geometry-Constrained Compiler}
\claim{把工具机理绑定到场景机理。}

这是方法核心。compiler 应被包装为从 public evidence 到 executable whole-tool poses 的桥。导师版只需讲约束族：contact alignment、motion-axis alignment、approach/clearance、stage-wise subgoals。输出是 fixed $SE(3)$ goals，不是低层动作，也不是 VLM trajectory。

\subsection{D. Fixed-Goal Dexterous Rollout and Evaluation}
\claim{让手执行，让 task contract 判断最终物理效果。}

policy 部分应保持克制：本文不声称训练了通用灵巧工具使用策略，而是使用 SimToolReal fixed-goal policy 作为执行层。评价层通过 task contracts 判断最终效果，并将失败归因到 grounding、compilation、execution 或 task effect。

\section{BENCHMARK AND EVALUATION PLAN}

\subsection{benchmark 应证明什么}
benchmark 要证明这个表述是 tool-centric 的，而不是 task-name-specific 的。理想状态下覆盖多类功能几何和至少两类场景，同时严格区分 headline statistics、candidate tasks、diagnostic tasks 和 draft/smoke tasks。

\subsection{任务分层}
\begin{table}[t]
\caption{导师讨论版任务分层。具体任务归属应在 contracts 和 launcher summaries 冻结后再定。}
\label{tab:tiers}
\begin{center}
\begin{tabular}{ll}
\toprule
层级 & 论文中的作用 \\
\midrule
Headline & 用于主结果表的冻结统计 \\
Candidate & 有潜力但仍需最终审查的任务 \\
Diagnostic & 用于诊断机理失败的压力测试 \\
Draft/Smoke & 开发证据，不作为最终 claim \\
\bottomrule
\end{tabular}
\end{center}
\end{table}

\subsection{评价指标}
指标要解释机理链条，而不只是报告成功率：functional grounding、scene grounding、subgoal validity、policy execution、final task success 和 first-failure attribution。没有 frozen launcher summaries 之前，不填具体数字。

\subsection{预期主结果表}
最终论文应有一个紧凑主表：
\begin{quote}
Task, Tool, Scene, $N$, grounding validity, subgoal validity, execution
completion, final success, dominant failure mode.
\end{quote}
具体列名等 contracts 稳定后再定。

\section{PAPER PACKAGING}

\subsection{一句话 pitch}
\begin{quote}
\method{} 不是让 VLM 编造机器人轨迹，而是把公开工具证据和场景证据编译成 tool-centric whole-tool subgoals，使灵巧工具使用变得可审计、可归因、可执行。
\end{quote}

\subsection{主图要传达的信息}
\begin{quote}
工具几何定义 functional frame；RGB-D 场景证据定义 effect frame；compiler 绑定二者；hand follows。
\end{quote}

\subsection{需要避免的过度 claim}
导师版应提醒自己：不要过度声称通用工具使用、通用灵巧性、在线自适应、real-robot transfer 或 open-ended VLM planning，除非实验支持。也不要把历史迁移快照当成最终成功率。

\section{OPEN DECISIONS FOR ADVISOR DISCUSSION}

\begin{itemize}
  \item 标题应该用当前 slogan 风格，还是更稳妥地突出 tool-centric visual grounding？
  \item 哪些任务在 contract 冻结后可以作为 headline statistics？
  \item RA-L 级别需要多少工具形态和场景形态才显得足够有说服力？
  \item 最公平的 baseline 是 direct VLM waypointing、去掉 tool axis、去掉 contact point，还是 primitive-style planning？
  \item 标题里是否要直接写 dexterous manipulation，还是把 dexterity 放在 execution 部分以避免过度承诺？
\end{itemize}

\section{LIMITATIONS}

局限性应作为机理边界的自然结果来写。如果功能部位不可见或未被充分 review，工具机理会错；如果深度 lifting 噪声大，场景机理会错；如果 compiler 给出几何可行的 tool pose，但 fixed-goal policy 无法保持接触，最终任务仍会失败。这些局限是可以接受的，因为本文没有用 repair loop 或 privileged simulator refs 把失败藏起来。

\section{CONCLUSION}

这篇文章最后应回到一个简单、可守的结论：工具使用是一个 mechanism-binding problem。\method{} 通过公开视觉证据显式恢复工具和场景机理，用几何编译器绑定二者，并交给 fixed-goal dexterous policy 执行。最强包装是：工具定义 frame，编译器让 frame 可执行，手负责跟随。

\bibliographystyle{IEEEtran}
\bibliography{references}

\end{document}
